% =================================================
% CHAPITRE 2 : VECTEURS
% =================================================

\chapter{Les vecteurs}

\section{Définition}

\begin{definition}[Vecteur]
Soit $n \in \N^*$. Un \emph{vecteur} de dimension $n$
est un $n$-uplet de réels, c'est-à-dire un élément
de $\R^n$ :
\[
    \vect{u} = \begin{pmatrix} u_1 \\ u_2 \\ \vdots \\ u_n \end{pmatrix},
    \qquad u_i \in \R \text{ pour tout } i \in \{1, \dots, n\}.
\]
Les réels $u_i$ sont les \emph{composantes} de
$\vect{u}$, et $n$ est sa \emph{dimension}
(son nombre de composantes).
\end{definition}

\begin{remark}
L'exposant de $\R^n$ est la dimension $n$, fixée
pour un vecteur donné. L'indice $i$, lui, parcourt
les composantes de $1$ à $n$.
\end{remark}

\begin{definition}[Vecteur nul]
Le vecteur nul de $\R^n$ est
$\vect{0} = (0, \dots, 0)^\top$ : toutes ses
composantes valent $0$.
\end{definition}

Dans toute la suite, $\vect{u}$ et $\vect{v}$ sont
deux vecteurs de $\R^n$ et $k \in \R$ est un réel.

\section{Opérations}

\subsection{Addition et soustraction}

\begin{definition}
Pour deux vecteurs \textbf{de même dimension} $n$ :
\[
    \vect{u} + \vect{v} =
    \begin{pmatrix} u_1 + v_1 \\ \vdots \\ u_n + v_n \end{pmatrix},
    \qquad
    \vect{u} - \vect{v} =
    \begin{pmatrix} u_1 - v_1 \\ \vdots \\ u_n - v_n \end{pmatrix}.
\]
\end{definition}

\begin{importantbox}
La somme d'un vecteur de $\R^n$ et d'un vecteur de
$\R^m$ avec $n \neq m$ n'est pas définie.
\end{importantbox}

\subsection{Multiplication par un réel}

\begin{definition}
\[
    k \vect{u} =
    \begin{pmatrix} k u_1 \\ \vdots \\ k u_n \end{pmatrix}.
\]
La division par un réel $k \neq 0$ est la
multiplication par $\frac{1}{k}$ :
$\dfrac{\vect{u}}{k} = \dfrac{1}{k}\,\vect{u}$.
\end{definition}

\subsection{Produit scalaire}

\begin{definition}[Forme analytique]
Pour deux vecteurs de même dimension :
\[
    \vect{u} \cdot \vect{v} = \sum_{i=1}^{n} u_i v_i
    = u_1 v_1 + u_2 v_2 + \dots + u_n v_n.
\]
Le résultat est un \textbf{réel}, pas un vecteur.
\end{definition}

\begin{proposition}[Forme géométrique]
Si $\vect{u}$ et $\vect{v}$ sont non nuls, et si
$\theta \in [0, \pi]$ est l'angle entre eux :
\[
    \vect{u} \cdot \vect{v}
    = \norm{\vect{u}}\,\norm{\vect{v}} \cos\theta.
\]
\end{proposition}

\begin{remark}
Les deux formes donnent le \emph{même} nombre. En
pratique, on calcule le produit scalaire par la forme
analytique, puis on s'en sert pour trouver l'angle :
\[
    \cos\theta =
    \frac{\vect{u} \cdot \vect{v}}{\norm{\vect{u}}\,\norm{\vect{v}}}
    \quad\text{donc}\quad
    \theta = \arccos\!\left(
    \frac{\vect{u} \cdot \vect{v}}{\norm{\vect{u}}\,\norm{\vect{v}}}
    \right).
\]
Si $\vect{u} \cdot \vect{v} = 0$, les vecteurs
sont \emph{orthogonaux} ($\theta = \frac{\pi}{2}$).
\end{remark}

\subsection{Norme}

\begin{definition}[Norme euclidienne]
\[
    \norm{\vect{u}} = \sqrt{\vect{u} \cdot \vect{u}}
    = \sqrt{\sum_{i=1}^{n} u_i^2}.
\]
C'est la longueur du vecteur. Elle est toujours
positive, et nulle seulement pour $\vect{u} = \vect{0}$.
\end{definition}

\subsection{Normalisation}

\begin{definition}
Pour $\vect{u} \neq \vect{0}$, le vecteur normalisé
(ou unitaire) de $\vect{u}$ est
\[
    \hat{\vect{u}} = \frac{1}{\norm{\vect{u}}}\,\vect{u}.
\]
Il a la même direction et le même sens que
$\vect{u}$, et sa norme vaut $1$.
\end{definition}

\begin{importantbox}
Le vecteur nul ne peut pas être normalisé : on
diviserait par $\norm{\vect{0}} = 0$.
\end{importantbox}

\section{Exemples calculés à la main}

Tous les exemples utilisent les deux vecteurs de $\R^3$
\[
    \vect{u} = \begin{pmatrix} 1 \\ 2 \\ 2 \end{pmatrix},
    \qquad
    \vect{v} = \begin{pmatrix} 2 \\ -2 \\ 1 \end{pmatrix}.
\]

\begin{example}[Addition et soustraction]
\[
    \vect{u} + \vect{v} =
    \begin{pmatrix} 1 + 2 \\ 2 + (-2) \\ 2 + 1 \end{pmatrix}
    = \begin{pmatrix} 3 \\ 0 \\ 3 \end{pmatrix},
    \qquad
    \vect{u} - \vect{v} =
    \begin{pmatrix} 1 - 2 \\ 2 - (-2) \\ 2 - 1 \end{pmatrix}
    = \begin{pmatrix} -1 \\ 4 \\ 1 \end{pmatrix}.
\]
\end{example}

\begin{example}[Multiplication et division par un réel]
Avec $k = 2$ :
\[
    2\vect{u} =
    \begin{pmatrix} 2 \times 1 \\ 2 \times 2 \\ 2 \times 2 \end{pmatrix}
    = \begin{pmatrix} 2 \\ 4 \\ 4 \end{pmatrix},
    \qquad
    \frac{\vect{u}}{2} =
    \begin{pmatrix} 0{,}5 \\ 1 \\ 1 \end{pmatrix}.
\]
\end{example}

\begin{example}[Produit scalaire analytique]
\[
    \vect{u} \cdot \vect{v}
    = 1 \times 2 + 2 \times (-2) + 2 \times 1
    = 2 - 4 + 2 = 0.
\]
Les vecteurs $\vect{u}$ et $\vect{v}$ sont donc
orthogonaux.
\end{example}

\begin{example}[Norme]
\[
    \norm{\vect{u}} = \sqrt{1^2 + 2^2 + 2^2}
    = \sqrt{9} = 3,
    \qquad
    \norm{\vect{v}} = \sqrt{2^2 + (-2)^2 + 1^2}
    = \sqrt{9} = 3.
\]
\end{example}

\begin{example}[Produit scalaire géométrique et angle]
Pour $\vect{u}$ et $\vect{v}$ :
\[
    \cos\theta = \frac{0}{3 \times 3} = 0
    \quad\text{donc}\quad
    \theta = \frac{\pi}{2},
\]
ce qui confirme l'orthogonalité trouvée plus haut.

Pour un cas où l'angle n'est pas droit, prenons
$\vect{a} = (1, 0)^\top$ et $\vect{b} = (1, 1)^\top$
dans $\R^2$ : $\vect{a} \cdot \vect{b} = 1$,
$\norm{\vect{a}} = 1$ et $\norm{\vect{b}} = \sqrt{2}$, donc
\[
    \cos\theta = \frac{1}{1 \times \sqrt{2}} = \frac{\sqrt{2}}{2}
    \quad\text{donc}\quad
    \theta = \frac{\pi}{4}.
\]
On vérifie la forme géométrique :
$\norm{\vect{a}}\,\norm{\vect{b}}\cos\theta
= 1 \times \sqrt{2} \times \frac{\sqrt{2}}{2} = 1
= \vect{a} \cdot \vect{b}$.
\end{example}

\begin{example}[Normalisation]
\[
    \hat{\vect{u}} = \frac{1}{3}\begin{pmatrix} 1 \\ 2 \\ 2 \end{pmatrix}
    = \begin{pmatrix} 1/3 \\ 2/3 \\ 2/3 \end{pmatrix},
    \qquad
    \norm{\hat{\vect{u}}}
    = \sqrt{\tfrac{1}{9} + \tfrac{4}{9} + \tfrac{4}{9}} = 1.
\]
\end{example}

\begin{example}[Cas d'erreur]
\begin{itemize}
    \item $(1, 2)^\top + (1, 2, 3)^\top$ n'est pas défini :
          les dimensions sont $2$ et $3$.
    \item $\vect{0} \in \R^3$ ne se normalise pas,
          car $\norm{\vect{0}} = 0$.
\end{itemize}
\end{example}

\section{Utilité dans le projet}

Un neurone reçoit des entrées $\vect{x} \in \R^n$,
possède un poids par entrée, rangés dans
$\vect{w} \in \R^n$, et un biais $b \in \R$. Sa sortie,
avant la fonction d'activation, est
\[
    z = \vect{w} \cdot \vect{x} + b.
\]
Le produit scalaire est donc le calcul de base de
tout le réseau. La condition « même dimension »
correspond à une exigence concrète : il faut autant de
poids que d'entrées.


\section{Implémentation}

\begin{codebox}
Voici un exemple d'implémentation en C++20.
\end{codebox}

\lstinputlisting[
language=C++,
caption=Code d'implémentation de la fonction membre "normalisation",
firstline=12,
lastline=23,
numbers=left,
breaklines=true
]{../../src/Vector.cpp}

\newpage
